\chapter{“世遗泉州”项目化学习方案设计}

\section{前期调查与需求分析}

在设计方案之前，研究者于2025年9月对所在班级的四十二名学生和六名语文教师进行了前期调查。学生问卷主要了解其对泉州世界遗产的熟悉程度、参与综合性学习的经验以及对项目化学习的期待；教师访谈则聚焦综合性学习的实施困难、资源使用习惯与评价方式。调查结果显示，学生对泉州世界遗产的总体了解较为有限，能够准确说出三个以上遗产点的学生不足三成；多数学生听说过“宋元泉州”“海丝文化”等说法，但对其历史背景与语文学习价值缺乏深入理解。与此同时，学生对走访、访谈、拍摄和展示表现出较高兴趣，希望综合性学习“不要只是查资料和做手抄报”。

教师访谈反映出三方面困难。第一，资源筛选困难。泉州世界遗产涉及历史、建筑、宗教、民俗等多个领域，教师难以在有限课时内判断哪些材料适合语文学习。第二，任务设计困难。综合性学习容易变成活动堆叠，教师缺少把活动转化为语文任务的经验。第三，评价困难。过程性表现难以量化，最终成果又容易由少数学生包办。基于这些调查结果，本研究将“资源筛选—任务转化—过程评价”作为方案设计需要重点解决的问题。

\section{设计原则}

结合前期调查与理论基础，本研究确定四条设计原则。第一，坚持文化性与语文性统一。项目可以围绕遗产展开，但每一项活动都必须指向具体的语文学习目标，例如提取信息、整理资料、撰写解说词、进行访谈或发表观点，避免把语文课上成历史文化课。第二，坚持真实性与可行性平衡。驱动性问题应来自真实情境，任务应具有可完成性，实地探究要综合考虑课时、安全、交通与家长支持等因素。第三，坚持任务链与支架并重。项目任务应形成由浅入深的链条，同时在关键节点提供阅读材料、观察记录表、访谈提纲、写作模板与评价量表，帮助学生完成挑战。第四，坚持过程性与多元评价。评价应覆盖资料搜集、合作讨论、实地探究、作品修改与展示反思等环节，采用学生自评、同伴互评与教师评价相结合的方式。

\section{项目主题与学习目标}

本项目的主题为“世遗泉州·青春讲解”，驱动性问题为：“如何让更多人读懂泉州世界遗产？”这一问题既指向真实的文化传播需求，也为学生的语言实践提供了对象与目的。项目的核心成果为一份面向中学生或游客的“世遗泉州”多模态导览作品，可以是图文手册、音频导览、短视频或校园展览，学生以小组为单位选择一种形式完成。成果必须包含遗产点介绍、现场观察记录、人物访谈、推荐路线与原创解说词，并在班级或校园公开展示。

依据《义务教育语文课程标准（2022年版）》与项目主题，本项目确定四类学习目标。文化理解方面，学生能够说出至少两处泉州世界遗产的基本信息，理解其历史背景与文化价值，并能结合具体材料表达自己的认识。语言运用方面，学生能够根据受众与目的选择恰当的表达方式，撰写准确、连贯、有吸引力的解说词，运用说明、描写、议论等多种表达方式。思维发展方面，学生能够围绕驱动性问题提出子问题，辨别事实与观点，整合多渠道信息，形成有依据的判断。合作探究方面，学生能够承担小组角色，参与讨论与分工，倾听他人意见，依据评价标准修改作品。

\section{项目任务与课时安排}

项目共安排十二课时，另设两次课外探究活动。整个项目分为入项、探究、创作与出项四个阶段。入项阶段通过情境视频、遗产图片与“如果你来当讲解员”的任务情境，激发学生的兴趣并明确项目成果；探究阶段围绕遗产点的历史信息、人物故事、语言材料与现场观察展开，引导学生完成资料卡片、访谈记录与观察笔记；创作阶段组织学生撰写解说词、设计导览路线、制作多模态作品，并通过同伴互评进行修改；出项阶段举办公开展示与反思活动，邀请其他班级学生、家长或社区文化工作者担任观众与评价者。具体任务与语文学习目标的对应关系如表~\ref{tab:project}所示。

\begin{table}[htbp]
  \centering
  \caption{“世遗泉州”项目任务与语文学习目标对应表}
  \label{tab:project}
  \begin{tabularx}{\textwidth}{@{}L{1.4cm}L{4.1cm}Yp{1.1cm}@{}}
    \toprule
    阶段 & 驱动任务 & 语文学习目标 & 课时 \\
    \midrule
    入项 & 观看遗产影像，提出“如何让更多人读懂泉州世遗”的子问题 & 提取关键信息，学习提出有价值的问题，明确成果的受众与目的 & 2 \\
    探究 & 走访遗产点，阅读地方文献，访谈讲解员或社区居民 & 筛选与整合信息，记录访谈要点，辨析事实与观点，积累语言材料 & 4 \\
    创作 & 撰写解说词，设计导览路线，制作手册、音频或短视频 & 根据受众调整表达，运用说明、描写、议论等方法，修改与润色作品 & 4 \\
    出项 & 举办“世遗泉州·青春讲解”展示会，开展自评与互评 & 当众表达与倾听，依据量表评价作品，撰写反思日志 & 2 \\
    \bottomrule
  \end{tabularx}
  \par\smallskip
  {\zihao{6}注：课时不含学生课外走访与资料整理时间。}
\end{table}

\section{项目实施流程}

项目实施流程如图~\ref{fig:flow}所示。第一阶段为“情境入项”，教师通过一段泉州世界遗产航拍视频和一组老照片，引导学生讨论“如果外地朋友来泉州，你会带他看什么”。学生在讨论中形成若干子问题，例如“洛阳桥为什么建在入海口”“浔埔女的服饰与海洋生活有什么关系”“开元寺的东西塔为什么能保存至今”。教师将这些子问题归类，帮助学生形成小组探究方向。

第二阶段为“任务分解与资料准备”。各小组选择一处遗产点，明确成果形式与受众，制定任务清单。教师提供“资料卡片模板”“访谈提纲模板”和“观察记录表”，并开设两节微型课，分别讲解如何筛选可靠资料、如何进行访谈记录。第三阶段为“实地探究与文本研读”。学生利用周末或综合实践课时走访遗产点，拍摄照片、记录现场信息，访谈讲解员、游客或社区居民；回到课堂后，教师组织学生阅读与遗产点相关的说明文、游记、新闻与地方文献，比较不同文本的表达方式，为撰写解说词积累材料。

第四阶段为“成果创作与迭代修改”。学生先完成解说词初稿，再依据“受众意识、信息准确、结构清晰、语言得体”四项标准进行同伴互评。教师针对共性问题开设写作微课，例如如何把专业术语转化为中学生能理解的语言，如何在解说中融入故事，如何用过渡句连接不同信息。第五阶段为“公开展示与反思”。各小组在校园展览区或班级展示会上呈现作品，观众通过贴纸投票与评价表给出反馈。学生根据反馈完成最终修改，并撰写反思日志，记录自己在资料搜集、合作分工、语言表达与问题解决方面的收获与不足。

\begin{figure}[htbp]
  \centering
  \begin{tikzpicture}[
    node distance=0.55cm and 0.55cm,
    box/.style={draw, rounded corners=3pt, align=center, minimum width=3.6cm,
      minimum height=1.1cm, font=\zihao{5}},
    arrow/.style={-{Stealth[length=2.5mm]}, thick}
  ]
    \node[box] (a) {情境入项\\明确驱动问题};
    \node[box, right=of a] (b) {任务分解\\确定成果与受众};
    \node[box, right=of b] (c) {实地探究\\观察与访谈};
    \node[box, below=of c] (d) {文本研读\\积累语言材料};
    \node[box, left=of d] (e) {成果创作\\撰写与修改};
    \node[box, left=of e] (f) {公开展示\\评价与反思};
    \draw[arrow] (a) -- (b);
    \draw[arrow] (b) -- (c);
    \draw[arrow] (c) -- (d);
    \draw[arrow] (d) -- (e);
    \draw[arrow] (e) -- (f);
    \draw[arrow] (f.south) -- ++(0,-0.35) -| node[pos=0.25, below, font=\zihao{6}]{迭代改进} (a.south);
  \end{tikzpicture}
  \caption{“世遗泉州”项目化学习实施流程}
  \label{fig:flow}
\end{figure}

\section{评价方案}

本项目采用表现性评价与过程性评价相结合的方式。评价对象包括资料卡片、访谈记录、观察笔记、解说词初稿与修改稿、最终作品、展示表现与反思日志。评价主体包括学生自评、同伴互评、教师评价与观众反馈。评价维度分为文化理解、语言运用、探究过程与合作表达四个方面，权重分别为百分之三十、百分之三十五、百分之二十与百分之十五。具体评价要点如表~\ref{tab:rubric}所示。

\begin{table}[htbp]
  \centering
  \caption{“世遗泉州”项目评价量表}
  \label{tab:rubric}
  \begin{tabularx}{\textwidth}{@{}L{2.2cm}L{4.2cm}Y@{}}
    \toprule
    评价维度 & 评价要点 & 表现水平描述 \\
    \midrule
    文化理解 & 能准确介绍遗产点信息，理解其历史背景与文化价值 & 能从多个角度解释遗产价值，并能联系当代生活表达自己的理解 \\
    语言运用 & 解说词准确、连贯，符合受众与目的 & 能根据展示对象调整语言，恰当运用说明、描写与议论，修改有针对性 \\
    探究过程 & 能提出子问题，筛选资料，完成观察与访谈 & 能辨别信息来源，记录关键细节，用证据支持自己的判断 \\
    合作表达 & 能承担小组角色，参与讨论与展示 & 能倾听并回应他人，分工合理，展示自然、有感染力 \\
    \bottomrule
  \end{tabularx}
\end{table}

为避免评价流于形式，研究者在项目开始前向学生公布评价量表，并用一份样例作品进行评分练习，帮助学生理解“优秀”“合格”“待改进”的具体表现。每次小组活动结束后，学生填写简短的过程记录，教师根据记录与课堂观察给予反馈。最终成绩由过程性证据与最终作品共同构成，其中过程性证据占百分之四十，最终作品占百分之六十。这样的安排意在引导学生重视探究过程，而不是只关注最后的展示效果。

在具体操作上，评价贯穿项目始终，而不是集中在最后一次展示。项目启动时，教师向学生说明评价维度与权重，并组织一次“样例評分”活动：师生共同阅读一份往届学生作品，依据量表逐项打分，讨论分数差异的原因。这一环节帮助学生把抽象的标准转化为具体的判断依据，也让学生意识到评价不是教师的“秘密裁决”，而是可以理解和参与的过程。项目进行中，每个阶段结束时设置一次简短的自评与互评，学生需要用“证据—解释—改进”的格式写下一条反馈，例如“我们组在解说词中加入了访谈对象的原话，这使内容更可信；但开头缺少对游客的称呼，建议增加问候语”。这样的反馈既指向作品，也指向合作过程。

评价量表也存在局限。文化理解中的情感认同、审美体验与价值判断，难以完全通过分数体现；合作中的倾听、妥协与责任担当，也不容易从最终作品中看出来。因此，本研究将量表评分与访谈、课堂观察、反思日志结合使用，避免用单一分数概括学生的学习经历。教师在使用量表时，应关注分数背后的证据，而不是把评价简化为“打分—排名”。对于表现暂时不理想的学生，应通过具体反馈指出改进方向，帮助其在后续任务中获得成功体验。

此外，学生也是评价规则的共同建构者。在项目开始前，教师可以邀请学生对评价量表提出修改建议，例如增加“作品是否适合目标受众”“是否尊重遗产地的文化语境”等条目。学生参与规则制定，有助于增强其对评价的认同感，也能促使学生从“被评价者”转变为“反思性实践者”。当学生能够用证据说明自己的进步、用标准解释他人的作品时，评价便不再只是教学的终点，而成为学习的一部分。

评价结果的使用同样需要谨慎。项目评价的目的不是给学生贴标签，而是为教学决策提供依据。教师可以通过过程记录发现小组合作的难点，通过作品分析判断学生在信息整合、结构安排或语言表达上的具体问题，并据此调整后续教学。例如，如果多数学生在解说词开头直接罗列信息，说明其受众意识尚未建立，教师可以在下一课时安排“开头方式比较”的微型课；如果学生在访谈记录中只写结论、不记原话，说明其证据意识不足，教师可以提供访谈录音转写的示范。评价只有转化为具体的教学行动，才能真正发挥促进作用。

需要指出的是，评价改革不能脱离学校整体制度。若学校仍以统一考试分数作为唯一评价依据，教师开展过程性评价的积极性就会受到影响。因此，学校层面应为项目化学习提供一定的制度空间，例如允许过程性评价计入平时成绩、减少重复性作业、为跨学科教研提供时间保障。只有当评价制度与课程改革相互支持，项目化学习才能持续开展，而不是成为一次性的展示活动。
